% Generated draft: requires mathematical and source review.
\documentclass[11pt]{article}
\usepackage[margin=25mm]{geometry}
\usepackage{fontspec}
\setmainfont{Latin Modern Roman}
\usepackage{amsmath,amssymb,mathtools,unicode-math}
\setmathfont{Latin Modern Math}
\usepackage{enumitem,needspace,xurl,hyperref}
\setlength{\parindent}{0pt}
\setlength{\parskip}{5pt}
\setlength{\emergencystretch}{3em}
\newcommand{\N}{\mathbb{N}}
\newcommand{\Z}{\mathbb{Z}}
\newcommand{\Q}{\mathbb{Q}}
\newcommand{\R}{\mathbb{R}}
\newcommand{\C}{\mathbb{C}}
\newcommand{\F}{\mathbb{F}}
\newcommand{\A}{\mathbb{A}}
\newcommand{\PP}{\mathbb{P}}
\newcommand{\E}{\mathbb{E}}
\newcommand{\eps}{\varepsilon}
\newcommand{\dd}{\,\mathrm{d}}
\begin{document}
\section*{UnsolvedMath ID: 8700007}
\textbf{Source problem code:} catalogue entry
\subsection{Definitions and mathematical statement}
Let $k\ge1$ and let $n_1<\cdots<n_k$ be positive integers. They are multiplicatively dependent if there exist integers $a_1,\ldots,a_k$, not all zero, such that
\[
 \prod_{i=1}^k n_i^{a_i}=1.
\]
This means joint dependence of the tuple, not necessarily pairwise dependence or minimal dependence of all its entries.

Langevin's conjecture as extracted is
\[
 \left(n_1<\cdots<n_k\text{ positive and multiplicatively dependent}\right)
 \quad\Longrightarrow\quad
 \exists p\text{ prime},\qquad n_1\le p\le n_k.
\]
Both interval endpoints are included. The assertion asks for a prime in the numerical interval, not necessarily a prime equal to an entry or dividing an entry.

The source does not explicitly require $k\ge2$ or $n_1>1$. Under the usual definition, $1$ creates multiplicative dependence, and the singleton tuple containing $1$ requires an exception to the literal prime-in-interval assertion. The intended exclusion or tuple-length convention must be checked rather than silently added. No further spacing, composite-entry, or minimal-dependence condition is supplied.
\subsection{Short English statement}
Must the interval from the smallest to the largest integer in a multiplicatively dependent increasing tuple contain a prime? The singleton-$1$ boundary case needs source qualification.

The supplied summary attributes the conjecture to Langevin as recorded by Waldschmidt, links it qualitatively to Grimm-type and Sylvester--Schur questions, and reports no located proof or disproof of the exact statement. The OPEN-TRIAGE label is a supplied assessment; provenance says a current primary-literature audit remains required.
\subsection{Sources}
\begin{itemize}
\item Local UnsolvedMath registry, permanent ID 8700007, source code catalogue entry, ``Conjecture 1.8 --- Langevin.'' Registry proposer metadata names Michel Waldschmidt; the title and summary attribute the conjecture itself to Langevin. Dataset curation and metadata: Ulam AI and UnsolvedMath Contributors, CC BY 4.0.
\item Michel Waldschmidt, \emph{Open Diophantine Problems} (2004), Conjecture 1.8, \url{https://arxiv.org/abs/math/0312440}; source-TeX extraction, accessed 2026-07-29.
\item Grimm-type and Sylvester--Schur problems are mentioned in the supplied summary without separate primary references.
\item Source-specific provenance: the survey presents the item as open; current primary-literature status audit and redistribution terms NEEDS\_REVIEW for the publicly accessible author/publisher copy, separately from the dataset licence.
\end{itemize}
\end{document}
